\documentclass[12pt]{article}
\usepackage[spanish,es-nodecimaldot]{babel}
\usepackage[utf8]{inputenc}\usepackage[T1]{fontenc}
\usepackage{newtxtext,newtxmath,microtype}
\usepackage[letterpaper,margin=2.55cm]{geometry}
\usepackage{setspace,csquotes}\onehalfspacing
\usepackage{amsmath,booktabs,tabularx,array,float,enumitem}
\usepackage{tikz}\usetikzlibrary{arrows.meta,positioning,shapes.geometric}
\usepackage{hyperref}\hypersetup{hidelinks,pdftitle={Contratos reconstruibles para sistemas NoSQL documentales},pdfauthor={Billy Jak Cordero Porras}}
\usepackage[backend=biber,style=apa,sorting=nyt,maxcitenames=2]{biblatex}
\addbibresource{referencias.bib}
\newcommand{\COBRIA}{\textsf{COBRIA}}
\title{\textbf{Contratos reconstruibles con ámbito explícito para sistemas NoSQL documentales}}
\author{Billy Jak Cordero Porras\\Investigador independiente, Coto Brus, Costa Rica}
\date{Agosto de 2026}

\begin{document}
\maketitle

\begin{abstract}
Los sistemas NoSQL documentales permiten adaptar la forma física de los datos a sus
patrones de acceso, pero esa flexibilidad puede producir copias derivadas sin autoridad,
aislamiento o reparación verificable. Este artículo formaliza el contrato de proyección
reconstruible con ámbito explícito de COBRIA. El contrato reúne fuente canónica,
transformación versionada, ámbito obligatorio, criterio de equivalencia, política de
frescura y procedimiento de reconstrucción. La propuesta se examina mediante ciencia
del diseño en tres perfiles documentales anonimizados y un protocolo de pruebas para
reconstrucción, idempotencia, evolución, recuperación autorizada y productos de datos.
La evidencia actual es estructural y ejecutable: demuestra aplicabilidad de los
contratos y permite formular hipótesis refutables. No se presentan cifras de rendimiento
obtenidas fuera de motores NoSQL ni se afirma superioridad universal. El resultado
principal es una especificación neutral y una agenda empírica reproducible para decidir
cuándo una proyección aporta valor sin convertirse en una segunda fuente de verdad.
\end{abstract}

\noindent\textbf{Palabras clave:} NoSQL documental; datos canónicos; proyecciones
reconstruibles; multiinquilinato; procedencia; minería de datos; inteligencia artificial.

\section{Introducción}

Una aplicación documental rara vez conserva una sola forma de cada hecho. Las reglas de
negocio necesitan documentos completos; las interfaces requieren respuestas breves;
la analítica necesita cortes estables; y la inteligencia artificial necesita fragmentos
con evidencia y autorización. Duplicar puede ser correcto, pero la copia se vuelve
peligrosa cuando nadie puede responder de dónde provino, qué versión representa, a qué
ámbito pertenece o cómo se repara.

La literatura describe separación de lectura y escritura, procesamiento de eventos,
procedencia, almacenes de características y gobierno de productos de datos
\parencite{microsoft2025cqrs,fowler2005eventsourcing,buneman2001provenance,
orr2021featurestores,li2023featurestore}. COBRIA no reclama esos mecanismos como nuevos.
Su contribución es composicional: convertir reconstrucción, ámbito y equivalencia en un
contrato verificable que acompaña a cada representación derivada.

La pregunta principal es: \emph{¿qué contrato permite optimizar lecturas documentales
sin crear una autoridad paralela ni debilitar aislamiento, trazabilidad y capacidad de
reparación?} Se derivan cuatro preguntas: cómo formalizar la proyección; qué invariantes
comprobar; cómo extender el contrato hacia analítica e inteligencia artificial; y qué
evidencia hace falta para justificar su costo.

Las contribuciones son:
\begin{enumerate}[leftmargin=*]
\item una definición formal del contrato reconstruible con ámbito;
\item una taxonomía neutral de canónicos, proyecciones, productos y mecanismos;
\item un protocolo NoSQL para lectura, escritura, reconstrucción, aislamiento y linaje;
\item una separación explícita entre evidencia disponible y evaluación pendiente.
\end{enumerate}

\section{Contexto y vacío de investigación}

Bigtable y Dynamo mostraron que partición, replicación, disponibilidad y patrones de
acceso deben considerarse desde el diseño \parencite{chang2006bigtable,
decandia2007dynamo}. La flexibilidad documental no elimina el esquema: lo distribuye
entre código, índices, reglas de acceso, eventos, cachés y consumidores.

La separación de modelos ofrece vocabulario para lectura y escritura; los eventos pueden
rehidratar estado; los almacenes de características preservan definiciones entre
entrenamiento e inferencia; y la procedencia explica cómo se obtuvo un resultado. Una
implementación todavía debe decidir quién posee autoridad, qué ámbito restringe la
copia, qué equivalencia se acepta y qué reparación es ejecutable.

El vacío abordado es la falta de una ficha contractual común que conecte optimización
operacional, aislamiento y productos para analítica e inteligencia artificial dentro de
sistemas documentales.

\section{Contrato COBRIA}

Sea $C$ una colección de documentos canónicos. Una proyección $P_i$ se describe por:
\[
K_i=\langle C,f_i,\theta_i,s,v,\phi_i,e_i,r_i,o_i\rangle,
\]
donde $f_i$ es la transformación, $\theta_i$ sus catálogos y políticas, $s$ el ámbito,
$v$ la versión, $\phi_i$ el límite de frescura, $e_i$ el verificador de equivalencia,
$r_i$ la reconstrucción y $o_i$ la propiedad operacional del componente.

El contrato satisface seis invariantes:
\begin{enumerate}[leftmargin=*]
\item \textbf{autoridad}: solo $C$ acepta cambios de dominio autoritativos;
\item \textbf{ámbito}: identidad, consulta, caché, evento y recuperación incluyen $s$;
\item \textbf{procedencia}: cada elemento derivado enlaza origen y versión;
\item \textbf{reconstrucción}: $r_i(C,\theta_i,v)$ produce una candidata completa;
\item \textbf{equivalencia}: $e_i$ compara resultado lógico, no representación física;
\item \textbf{operabilidad}: existen responsable, métricas, alerta y criterio de retiro.
\end{enumerate}

La publicación segura construye $P_i^{v+1}$, verifica equivalencia y cambia un alias de
forma indivisible. La versión activa no se vacía durante el proceso. La actualización
incremental es válida si puede repetirse y si una reconstrucción total periódica detecta
omisiones o deriva.

\section{Arquitectura documental}

\begin{figure}[H]\centering
\begin{tikzpicture}[node distance=1.1cm and 1.25cm,b/.style={draw,rounded corners,minimum width=3.1cm,minimum height=.9cm,align=center},a/.style={-{Latex[length=2mm]},thick}]
\node[b] (u) {Caso de uso};
\node[b,below left=of u] (c) {Documento canónico};
\node[b,below=of u] (p) {Proyección reconstruible};
\node[b,below right=of u] (d) {Producto de datos};
\node[b,below=of p] (m) {Metadatos, ámbito y versión};
\draw[a] (u)--(c);\draw[a] (c)--(p);\draw[a] (p)--(u);
\draw[a] (p)--(d);\draw[a] (m)--(c);\draw[a] (m)--(p);\draw[a] (m)--(d);
\end{tikzpicture}
\caption{Relaciones principales del contrato COBRIA en NoSQL documental.}
\end{figure}

El objeto canónico expresa identidad, revisión, tiempo y ciclo de vida. El Dater traduce
entre esa semántica y documentos físicos, índices y consultas. El servicio ejecuta
reglas y autorización. El catálogo conserva vocabulario versionado. Esta separación es
una guía de responsabilidades; no obliga a multiplicar clases cuando un módulo pequeño
puede cumplir el mismo contrato.

\section{Minería de datos e inteligencia artificial}

Un conjunto de datos reconstruible es una proyección fechada. Declara población,
ventana temporal, transformación, versión de catálogos, exclusiones, semilla y política
de división. Una característica versionada añade definición, tipo, tiempo de evento,
tiempo de conocimiento y prevención de fuga temporal.

La recuperación aumentada restringe candidatos por ámbito antes de calcular similitud.
Cada fragmento conserva documento fuente, revisión, versión, autorización y fecha. La
respuesta enlaza evidencia o se abstiene. Una herramienta de agente añade capacidad
mínima, validación, idempotencia, confirmación para efectos de alto riesgo y registro.

Estas propiedades no garantizan exactitud de un modelo. Sí hacen reproducible el insumo,
limitan cruces de ámbito y permiten retirar productos afectados por cambios o borrado.

\section{Método}

La investigación sigue ciencia del diseño: problema, objetivos, construcción,
demostración, evaluación y comunicación. Tres perfiles anonimizados representan eventos
y multimedia, operación hotelera y operación comercial. Se inspeccionan canónicos,
Daters, servicios, catálogos, reglas, índices, funciones y derivados.

La evaluación activa se restringe a NoSQL documental. Se distinguen evidencia
estructural, ejecutable y empírica. El protocolo registra motor, región, escala, ámbitos,
forma documental, índices, semillas, calentamiento, orden, versión y descartes. La
corrida completa es la unidad inferencial.

Las hipótesis son: HN1, reconstrucción equivalente; HN2, cero cruces de ámbito; HN3,
idempotencia y revisión; HN4, reproducción de productos analíticos; y HN5, beneficio de
lectura suficiente para pagar mantenimiento. HN5 permanece abierta hasta ejecutar una
réplica multiescala en motores documentales.

\section{Protocolo experimental NoSQL}

\begin{table}[H]\centering
\caption{Experimentos y criterios de aceptación.}
\begin{tabularx}{\textwidth}{p{1.4cm}p{3.5cm}X}\toprule
Código & Experimento & Aceptación\\\midrule
P1 & consulta canónica frente a proyección & mismo resultado y mejora del objetivo declarado\\
P2 & amplificación de escritura & operaciones, bytes, eventos y convergencia dentro de presupuesto\\
R1 & reconstrucción total e incremental & equivalencia, idempotencia y publicación segura\\
S1 & aislamiento adversarial & cero documentos cruzados y rechazo sin ámbito\\
A1 & conjunto, recuperación y herramientas & procedencia completa, abstención y cero fugas\\
\bottomrule\end{tabularx}\end{table}

P1 y P2 deben ejecutarse en al menos dos motores NoSQL, con treinta corridas válidas por
celda, orden rotado, escalas múltiples y resultados negativos publicados. R1 interrumpe
la construcción antes de publicar; S1 manipula identidad, caché y búsqueda; A1 introduce
eventos tardíos, duplicados, catálogos desconocidos y documentos incompletos.

\section{Resultados disponibles}

Los tres perfiles contienen fuentes documentales, capas de acceso, servicios, catálogos
y derivados. La recurrencia confirma que el problema existe en dominios heterogéneos.
Las pruebas disponibles verifican reglas de acceso, roles, aislamiento, ciclos de vida,
idempotencia y reconstrucción en rutas concretas. Una observación longitudinal también
muestra reducción de accesos directos desde presentación hacia la capa documental.

Esta evidencia apoya viabilidad estructural y ejecutable. No demuestra beneficio
universal de rendimiento, productividad humana ni seguridad exhaustiva. Las cifras de
pilotos ajenos al alcance NoSQL se excluyeron y no alimentan el resumen ni las
conclusiones.

\section{Proposiciones verificables}

El contrato permite derivar propiedades que un arnés puede intentar refutar. No se
presentan como teoremas universales del almacenamiento distribuido, sino como
obligaciones de una implementación COBRIA.

\subsection{Autoridad única}

Sea $w(x)$ una operación que cambia el estado de dominio de $x$. Para cualquier
proyección $P_i$, una operación válida satisface:
\[
w(x)\Rightarrow w(C_x)\land \neg w_{autoridad}(P_i(x)).
\]
Una interfaz administrativa que modifica directamente el derivado viola la propiedad,
aunque luego intente sincronizar la fuente. La prueba combina análisis estático de
dependencias con una escritura adversarial contra el adaptador.

\subsection{Aislamiento previo a recuperación}

Para una identidad $a$ y ámbito autorizado $s_a$, el conjunto candidato cumple:
\[
Cand(a,q)\subseteq \{x\mid scope(x)=s_a\land permit(a,x)\}.
\]
La condición se aplica antes de similitud, orden, caché o ensamblado de contexto. Así,
un documento prohibido no influye en puntuaciones, tiempos, conteos ni explicaciones.

\subsection{Reconstrucción equivalente}

Para una versión congelada $v$, política $\theta$ y corte $t$, la reconstrucción se
acepta cuando:
\[
e_i\bigl(P_i^{activo},r_i(C_{\leq t},\theta,v)\bigr)=1.
\]
El verificador puede comparar igualdad exacta, conjunto ordenado, tolerancia numérica o
invariantes de dominio. La tolerancia debe declararse antes de observar divergencias.

\subsection{Monotonicidad de revisión}

Si $rev(x')<rev(C_x)$, el procesamiento de $x'$ no puede reducir la revisión visible ni
reaplicar un efecto no conmutativo. La propiedad cubre entregas fuera de orden,
reintentos y restauraciones parciales.

\section{Modelo de decisión multidimensional}

Una proyección no se justifica únicamente por latencia. Se define el vector:
\[
D_i=\langle \Delta L,\Delta O,\Delta B,\Delta S,\Delta F,\Delta R,\Delta H\rangle,
\]
donde $L$ es latencia, $O$ operaciones facturables, $B$ bytes, $S$ almacenamiento,
$F$ frescura, $R$ riesgo operacional y $H$ esfuerzo humano. Los tres primeros pueden
medirse automáticamente; los últimos requieren objetivos, incidentes y estudios humanos.

La frontera temporal simplificada usa consultas $Q$, escrituras $W$, ahorro de lectura
$\Delta q$ y costo adicional $\Delta w$:
\[
Q\Delta q>W\Delta w.
\]
Esta desigualdad no decide por sí sola. Si una proyección duplica información sensible,
incumple frescura o carece de reparación, se rechaza aunque el tiempo agregado resulte
favorable. El valor del modelo es obligar a declarar dimensiones y evitar que una
mediana aislada se convierta en recomendación arquitectónica.

\begin{table}[H]\centering
\caption{Preguntas de decisión antes de adoptar una proyección.}
\begin{tabularx}{\textwidth}{p{3.3cm}X X}\toprule
Dimensión & Evidencia mínima & Señal para simplificar\\\midrule
Lectura & percentiles pareados y costo & índice canónico cumple el objetivo\\
Escritura & operaciones, eventos y convergencia & amplificación supera presupuesto\\
Frescura & atraso y distribución temporal & consumidor exige lectura inmediata\\
Reparación & ensayo total e incremental & reconstrucción no es repetible\\
Privacidad & pruebas negativas por ámbito & derivado duplica campos innecesarios\\
Operación & responsable, alertas e incidentes & nadie posee o utiliza la proyección\\
\bottomrule\end{tabularx}\end{table}

\section{Ejemplo de aplicación neutral}

Considérese una colección de actividades con documentos heterogéneos, diez ámbitos y una
pantalla que muestra únicamente identidad, fecha, estado y resumen. La fuente canónica
conserva además reglas, multimedia, revisión y procedencia. El equipo observa que la
consulta indexada cumple el objetivo en escalas pequeñas, pero aumenta operaciones y
bytes en la escala superior.

La ficha COBRIA no ordena crear una copia. Primero define la consulta y congela el índice
canónico. Después propone una proyección estrecha con los cuatro campos visibles,
revisión, fuente y versión. P1 compara ambas rutas; P2 mide la actualización; R1 elimina
y reconstruye la candidata; S1 intenta cruzar ámbitos; A1 verifica que un conjunto
analítico pueda enlazar cada registro con el documento de origen.

Tres resultados son posibles. Si la consulta canónica satisface costo y latencia, se
descarta la proyección. Si la proyección mejora lectura pero rompe frescura o reparación,
se corrige antes de publicarla. Solo si cumple el vector de decisión se adopta, con fecha
de revisión y criterio de retiro. El ejemplo muestra que COBRIA es un proceso de
decisión, no una instrucción de duplicar.

\section{Trazabilidad entre hipótesis y artefactos}

\begin{table}[H]\centering
\caption{Matriz de trazabilidad del artículo.}
\begin{tabularx}{\textwidth}{p{1.5cm}p{3.3cm}p{3.2cm}X}\toprule
Hipótesis & Riesgo & Prueba & Evidencia requerida\\\midrule
HN1 & autoridad paralela & R1 y escritura adversarial & equivalencia y rechazo directo\\
HN2 & fuga entre ámbitos & S1 & cero candidatos cruzados\\
HN3 & duplicados o regresión & repetición y revisión vieja & un efecto y revisión monotónica\\
HN4 & producto irreproducible & A1 & manifiesto, huella y procedencia\\
HN5 & costo sin beneficio & P1 y P2 & vector pareado por región de carga\\
\bottomrule\end{tabularx}\end{table}

La matriz evita que una métrica secundaria sustituya la hipótesis. Por ejemplo, una
reconstrucción rápida no apoya HN1 si produce documentos distintos; una precisión alta
no apoya HN2 si el conjunto candidato contenía material prohibido.

\section{Plan de análisis y ciencia abierta}

Cada celda contiene treinta corridas válidas. Se publican medianas, intervalos del 95\%
obtenidos mediante remuestreo de corridas, diferencias y razones pareadas. La dirección
de cada par se conserva. Las familias confirmatorias se ajustan por comparaciones
múltiples y las exploraciones se rotulan por separado.

El paquete de réplica debe incluir protocolo, enmiendas, generador, adaptadores,
manifiestos, eventos crudos, análisis, tablas y entorno. Las huellas enlazan cada figura
con sus entradas. Una corrida se descarta únicamente por reglas congeladas: manifiesto
incompleto, pérdida de equivalencia previa a la medición, interferencia externa registrada
o fallo del arnés. Los datos descartados se conservan fuera del conjunto inferencial.

La revisión independiente examinará equivalencia de configuraciones, unidad inferencial,
exclusiones y correspondencia entre evidencia y afirmaciones. El artículo se actualizará
con cifras solo después de completar dos motores documentales y reproducir las salidas
desde el depósito.

\section{Discusión}

La reconstruibilidad cambia la naturaleza de una copia: deja de ser un activo que debe
protegerse manualmente y se convierte en un producto regenerable con contrato. El ámbito
explícito evita que optimización y recuperación debiliten la autorización. La
procedencia continua conecta operación, analítica e inteligencia artificial.

El costo es real: más metadatos, observabilidad, almacenamiento, eventos y disciplina.
Por ello COBRIA incorpora un criterio de retiro. Una proyección sin objetivo,
responsable, verificador o uso suficiente debe eliminarse. En sistemas pequeños, la
mejor decisión puede ser conservar solo documentos canónicos bien indexados.

\section{Amenazas a la validez}

Los perfiles comparten autoría y ecosistema. Los emuladores no reproducen todas las
cuotas, particiones, regiones o costos administrados. La inspección no sustituye
ejecución; los datos sintéticos pueden omitir errores históricos; y las pruebas
automáticas no prueban comprensión humana. Se requieren réplica externa, revisión
independiente y estudio con desarrolladores ajenos a los casos.

\section{Conclusión}

COBRIA formaliza una arquitectura NoSQL documental en la que cada derivado conserva
autoridad rastreable, ámbito, versión, equivalencia y reconstrucción. Su aporte no es
inventar repositorios, eventos o proyecciones, sino integrarlos en una especificación
verificable que alcanza productos analíticos y de inteligencia artificial.

La evidencia actual permite afirmar aplicabilidad y comprobación de contratos en rutas
concretas. La afirmación cuantitativa fuerte permanece abierta hasta completar P1--A1 en
motores NoSQL y publicar el paquete congelado. Esta delimitación transforma una intuición
arquitectónica en un programa científico refutable.

\printbibliography
\end{document}
